% Évaluation pratique — Tableur LibreOffice Calc (BTS CIEL 1re année) — 45 min
% Compilation :
%   sujet   : pdflatex Eval_TP_Tableur.tex
%   corrigé : pdflatex -jobname=Eval_TP_Tableur_Corrige "\def\modecorrige{1}% Évaluation pratique — Tableur LibreOffice Calc (BTS CIEL 1re année) — 45 min
% Compilation :
%   sujet   : pdflatex Eval_TP_Tableur.tex
%   corrigé : pdflatex -jobname=Eval_TP_Tableur_Corrige "\def\modecorrige{1}% Évaluation pratique — Tableur LibreOffice Calc (BTS CIEL 1re année) — 45 min
% Compilation :
%   sujet   : pdflatex Eval_TP_Tableur.tex
%   corrigé : pdflatex -jobname=Eval_TP_Tableur_Corrige "\def\modecorrige{1}% Évaluation pratique — Tableur LibreOffice Calc (BTS CIEL 1re année) — 45 min
% Compilation :
%   sujet   : pdflatex Eval_TP_Tableur.tex
%   corrigé : pdflatex -jobname=Eval_TP_Tableur_Corrige "\def\modecorrige{1}\input{Eval_TP_Tableur.tex}"
\documentclass[10pt,a4paper]{article}

\usepackage[T1]{fontenc}
\usepackage[utf8]{inputenc}
\usepackage{lmodern}
\usepackage[french]{babel}
\usepackage[a4paper,margin=1.5cm,top=1.3cm,bottom=1.8cm]{geometry}
\usepackage[table]{xcolor}
\usepackage{booktabs,tabularx,array}
\usepackage{enumitem}
\usepackage{fancyhdr}
\usepackage{amssymb}
\usepackage[most]{tcolorbox}
\usepackage{hyperref}
\hypersetup{hidelinks}

\providecommand{\modecorrige}{0}
\newif\ifcorrige
\ifnum\modecorrige=1 \corrigetrue \else \corrigefalse \fi

\definecolor{ipred}{HTML}{C0392B}
\definecolor{cdcbar}{HTML}{A67C3D}
\definecolor{cdcbg}{HTML}{F4F1EA}
\definecolor{qbar}{HTML}{4A6B5D}
\definecolor{qbg}{HTML}{EDF1EF}
\definecolor{corr}{HTML}{1F5FA8}
\definecolor{corrbg}{HTML}{E8F0FA}
\definecolor{grisfoot}{HTML}{828282}
\definecolor{entete}{HTML}{D9E2DD}

\newcommand{\cel}[1]{\textcolor{ipred}{\texttt{#1}}}   % référence de cellule / formule
\newcommand{\pts}[1]{\hfill\mbox{\textbf{(#1)}}}
\newcommand{\menu}[1]{\textsf{\small #1}}               % chemin de menu LibreOffice

\newtcolorbox{cdc}{enhanced,breakable,colback=cdcbg,colframe=cdcbar,boxrule=0pt,
  leftrule=3pt,arc=0pt,outer arc=0pt,left=5pt,right=5pt,top=2pt,bottom=2pt}
\newtcolorbox{partie}[2]{enhanced,breakable,colback=qbg,colframe=qbar,boxrule=0pt,
  leftrule=3pt,arc=0pt,outer arc=0pt,left=5pt,right=5pt,top=2pt,bottom=2pt,
  title={#1\hfill(#2)},coltitle=qbar,fonttitle=\bfseries,detach title,before upper={\tcbtitle\par\smallskip}}

% Zone de réponse : vierge dans le sujet, remplie dans le corrigé
\newcommand{\zone}[2]{%
  \ifcorrige
    \begin{tcolorbox}[colback=corrbg,colframe=corr,boxrule=0.6pt,arc=1pt,left=4pt,right=4pt,top=1pt,bottom=1pt,
      fontupper=\small\color{corr}]#2\end{tcolorbox}%
  \else
    \begin{tcolorbox}[colback=white,colframe=black!45,boxrule=0.5pt,arc=1pt,height=#1]\end{tcolorbox}%
  \fi}

\pagestyle{fancy}
\fancyhf{}
\renewcommand{\headrulewidth}{0pt}
\fancyfoot[L]{\small\itshape\color{grisfoot}Évaluation — Tableur LibreOffice Calc\ifcorrige\ — CORRIGÉ\fi}
\fancyfoot[R]{\small\color{grisfoot}\thepage}

\setlist{itemsep=0pt,topsep=1pt,parsep=0pt,leftmargin=*}
\setlength{\parindent}{0pt}
\setlength{\parskip}{3pt}

\begin{document}

\begin{center}
  {\Large\bfseries Évaluation pratique — Tableur (LibreOffice Calc)\par}
  \vspace{1pt}
  {\bfseries Formules, fonctions, mise en forme, tri et graphique — durée : 45 minutes\par}
  \vspace{2pt}
  BTS CIEL 1re année
  \ifcorrige\par\vspace{2pt}{\large\bfseries\color{corr}CORRIGÉ ET BARÈME — document enseignant}\fi
\end{center}

\vspace{-2pt}
{\renewcommand{\arraystretch}{1.5}
\begin{tabularx}{\linewidth}{|X|X|p{2.2cm}|}
\hline
\textbf{Nom :} & \textbf{Prénom :} & \textbf{Note :} \hfill /20 \\ \hline
\end{tabularx}}

\begin{cdc}
\textbf{Consignes.} Travail individuel, sans Internet. Documents de cours autorisés.
\begin{itemize}
  \item Ouvrir le fichier \texttt{Eval\_Tableur\_donnees.ods} (ou, à défaut, saisir le tableau ci-dessous en respectant
        \textbf{exactement} les cellules), puis l'\textbf{enregistrer immédiatement} sous
        \texttt{NOM\_Prenom\_Eval\_Tableur.ods} (\menu{Fichier > Enregistrer sous}). Enregistrer régulièrement.
  \item \textbf{Toute valeur calculée doit l'être par une formule} : un nombre recopié à la main vaut 0.
  \item Les réponses de la partie D se rédigent sur cette feuille, à rendre avec le fichier.
\end{itemize}
\end{cdc}

\textbf{Contexte.} L'entreprise \emph{NetCIEL} vend du matériel réseau. Le responsable commercial vous confie le
tableau de ses ventes 2025 (quantités vendues par trimestre) pour établir un bilan.

\begin{center}\small
\setlength{\tabcolsep}{4pt}
\begin{tabular}{|>{\columncolor{black!6}}c|l|c|c|c|c|c|c|c|c|c|}
\hline
\rowcolor{black!6} & \textbf{A} & \textbf{B} & \textbf{C} & \textbf{D} & \textbf{E} & \textbf{F} & \textbf{G} & \textbf{H} & \textbf{I} & \textbf{J} \\ \hline
1 & \multicolumn{10}{l|}{\textbf{Ventes 2025 — NetCIEL}} \\ \hline
3 & Produit & Prix unitaire & T1 & T2 & T3 & T4 & Total qté & Moy./trim. & CA annuel & Part du CA \\ \hline
4 & Switch 24 ports & 180 & 12 & 15 & 9 & 14 & & & & \\ \hline
5 & Routeur & 250 & 8 & 6 & 10 & 7 & & & & \\ \hline
6 & Point d'accès Wi-Fi & 95 & 20 & 25 & 18 & 30 & & & & \\ \hline
7 & Pare-feu & 420 & 3 & 5 & 4 & 6 & & & & \\ \hline
8 & Câble RJ45 Cat6 (lot de 10) & 35 & 40 & 55 & 38 & 60 & & & & \\ \hline
9 & Onduleur & 150 & 5 & 4 & 7 & 6 & & & & \\ \hline
10 & Total & & & & & & & & & \\ \hline
\end{tabular}
\end{center}
\vspace{-4pt}
{\footnotesize La ligne 2 reste vide. Le prix unitaire est en euros hors taxes ; CA = chiffre d'affaires.}

%---------------------------------------------------------------------
\begin{partie}{Partie A — Calculs avec formules et recopie}{8 pts}
\begin{enumerate}[label=\textbf{A\arabic*.}]
  \item En \cel{G4}, calculer le nombre total d'articles vendus dans l'année avec la fonction \cel{SOMME},
        puis recopier la formule jusqu'en \cel{G9} (poignée de recopie). \pts{1,5 pt}
  \item En \cel{H4}, calculer la moyenne trimestrielle des ventes avec la fonction \cel{MOYENNE}, puis recopier
        jusqu'en \cel{H9}. \pts{1 pt}
  \item En \cel{I4}, calculer le chiffre d'affaires annuel du produit (prix unitaire $\times$ total des quantités),
        puis recopier jusqu'en \cel{I9}. \pts{1,5 pt}
  \item Ligne 10 : calculer en \cel{C10} le total des ventes du 1\textsuperscript{er} trimestre, puis recopier
        la formule vers la droite jusqu'en \cel{G10}. Calculer aussi le CA total en \cel{I10}. \pts{1,5 pt}
  \item En \cel{J4}, calculer la part du CA du produit dans le CA total (\cel{I4} divisé par le CA total).
        La formule doit pouvoir être recopiée jusqu'en \cel{J9} \textbf{sans être modifiée} :
        utiliser une \textbf{référence absolue} (symbole \cel{\$}). \pts{1,5 pt}
  \item Saisir en \cel{A12} « CA le plus élevé », en \cel{A13} « CA le plus faible », en \cel{A14} « Écart ».
        Calculer en \cel{B12} et \cel{B13} ces valeurs à partir de la plage \cel{I4:I9} avec les fonctions
        \cel{MAX} et \cel{MIN}, puis leur écart en \cel{B14}. \pts{1 pt}
\end{enumerate}
\end{partie}

%---------------------------------------------------------------------
\begin{partie}{Partie B — Mise en forme}{4 pts}
\begin{enumerate}[label=\textbf{B\arabic*.}]
  \item Fusionner les cellules \cel{A1:J1} ; titre centré, en gras, taille 14. \pts{1 pt}
  \item En-têtes de la ligne 3 en gras, centrés, avec une couleur d'arrière-plan ; ajuster la largeur des colonnes
        pour que tout soit lisible. \pts{1 pt}
  \item Ajouter toutes les bordures au tableau \cel{A3:J10}. \pts{0,5 pt}
  \item Format \textbf{monétaire en euros} pour les colonnes \cel{B} et \cel{I} et pour \cel{B12:B14} ;
        format \textbf{pourcentage avec 1 décimale} pour la colonne \cel{J}. \pts{1,5 pt}
\end{enumerate}
\end{partie}

\newpage
%---------------------------------------------------------------------
\begin{partie}{Partie C — Tri et graphique}{5 pts}
\begin{enumerate}[label=\textbf{C\arabic*.}]
  \item Trier les produits (lignes 4 à 9 seulement, \textbf{sans la ligne Total}) par CA annuel
        \textbf{décroissant} (\menu{Données > Trier}). Vérifier que les résultats sont restés justes. \pts{1,5 pt}
  \item Insérer un graphique en \textbf{secteurs} représentant la répartition du CA annuel par produit
        (noms des produits et CA ; touche \textbf{Ctrl} pour sélectionner deux plages non adjacentes)
        (\menu{Insertion > Diagramme}). \pts{2 pts}
  \item Donner au graphique le titre « Répartition du CA 2025 » et afficher les pourcentages sur les secteurs.
        Le placer à côté du tableau, sans masquer de données. \pts{1,5 pt}
\end{enumerate}
\end{partie}

%---------------------------------------------------------------------
\begin{partie}{Partie D — Analyse (à rédiger sur cette feuille)}{3 pts}
\begin{enumerate}[label=\textbf{D\arabic*.}]
  \item Écrire la formule saisie en \cel{J4}. Que se serait-il passé en \cel{J5} si l'on avait recopié la formule
        \cel{=I4/I10}, sans le symbole \cel{\$} ? Expliquer. \pts{1 pt}
\zone{2.6cm}{\cel{=I4/\$I\$10}. Sans \cel{\$}, la recopie décale les deux références : \cel{J5} contiendrait
\cel{=I5/I11}. \cel{I11} est vide, d'où l'erreur \cel{\#DIV/0!}. Le \cel{\$} bloque la colonne et la ligne de la cellule du total.}
  \item Le produit le plus vendu en quantité est-il celui qui rapporte le plus ? Justifier avec vos résultats. \pts{1 pt}
\zone{2.4cm}{Non. Le câble RJ45 est le plus vendu (193 lots) mais ne rapporte que 6\,755~€ (15,6~\%) à cause de son
faible prix (35~€). Le switch rapporte le plus (9\,000~€, 20,8~\%) avec 50 unités seulement.}
  \item Le prix du routeur passe à 270~€. Quelles cellules sont modifiées automatiquement ? Pourquoi ? \pts{1 pt}
\zone{2.4cm}{Le CA du routeur (colonne I), le CA total \cel{I10}, toutes les parts de la colonne J, éventuellement
\cel{B12:B14}, et le graphique. Les formules utilisent des références de cellules : tout est recalculé dès qu'une donnée change.}
\end{enumerate}
\end{partie}

\textbf{Bonus (+1 pt)} — En \cel{K4}, afficher « Best-seller » si le total des quantités est supérieur ou égal
à 50, et rien sinon (fonction \cel{SI}), puis recopier jusqu'en \cel{K9}.

\textbf{En fin d'épreuve} : enregistrer, fermer le fichier, le déposer dans le dossier indiqué, rendre cette feuille.

%=====================================================================
\ifcorrige
\newpage
\section*{Corrigé — formules et résultats attendus}

LibreOffice Calc en français : noms de fonctions en français, séparateur d'arguments \cel{;}.
Valeurs ci-dessous \textbf{avant le tri} (après le tri, l'ordre des lignes change mais les valeurs restent les mêmes).

\begin{center}\small
\setlength{\tabcolsep}{4pt}
\begin{tabular}{@{}lcccc@{}}
\toprule
\textbf{Produit} & \textbf{G : Total qté} & \textbf{H : Moy./trim.} & \textbf{I : CA annuel} & \textbf{J : Part du CA} \\
\midrule
Switch 24 ports              & 50  & 12,5  & 9\,000,00~€ & 20,8~\% \\
Routeur                      & 31  & 7,75  & 7\,750,00~€ & 17,9~\% \\
Point d'accès Wi-Fi          & 93  & 23,25 & 8\,835,00~€ & 20,5~\% \\
Pare-feu                     & 18  & 4,5   & 7\,560,00~€ & 17,5~\% \\
Câble RJ45 Cat6 (lot de 10)  & 193 & 48,25 & 6\,755,00~€ & 15,6~\% \\
Onduleur                     & 22  & 5,5   & 3\,300,00~€ & 7,6~\% \\
\midrule
\textbf{Total (ligne 10)} & \textbf{407} & — & \textbf{43\,200,00~€} & (100~\%) \\
\bottomrule
\end{tabular}

\smallskip
Totaux trimestriels : \cel{C10} = 88 ; \cel{D10} = 110 ; \cel{E10} = 86 ; \cel{F10} = 123.
\quad \cel{B12} = 9\,000~€ ; \cel{B13} = 3\,300~€ ; \cel{B14} = 5\,700~€.
\end{center}

\begin{center}\small
\begin{tabularx}{\linewidth}{@{}lp{4.6cm}X@{}}
\toprule
\textbf{Question} & \textbf{Formule attendue} & \textbf{Barème et points de vigilance} \\
\midrule
A1 & \cel{=SOMME(C4:F4)} & 1 pt formule, 0,5 pt recopie. \cel{=C4+D4+E4+F4} : 1 pt sur 1,5 (fonction demandée). \\
A2 & \cel{=MOYENNE(C4:F4)} & \cel{=G4/4} accepté pour 0,5 pt sur 1 (fonction demandée). \\
A3 & \cel{=B4*G4} & 1 pt formule, 0,5 pt recopie. \\
A4 & \cel{=SOMME(C4:C9)} ; \cel{I10} : \cel{=SOMME(I4:I9)} & 1 pt ligne C–G recopiée, 0,5 pt \cel{I10}. \cel{H10} vide ou moyenne : accepté. \\
A5 & \cel{=I4/\$I\$10} & \cel{=I4/I\$10} accepté. Sans \cel{\$} : \cel{\#DIV/0!} dès \cel{J5}, 0,5 pt max. \\
A6 & \cel{=MAX(I4:I9)} ; \cel{=MIN(I4:I9)} ; \cel{=B12-B13} & 0,5 pt MAX/MIN, 0,5 pt écart. Valeur tapée : 0. \\
B1--B4 & — & Contrôle visuel. B4 : 1 pt monétaire, 0,5 pt pourcentage à 1 décimale. \\
C1 & — & Ordre attendu : Switch, Point d'accès, Routeur, Pare-feu, Câble, Onduleur. Si la ligne Total est incluse
dans le tri ou si seule la colonne I a été triée (données mélangées) : 0. \\
C2--C3 & — & Plages \cel{A4:A9} et \cel{I4:I9}. Graphique en secteurs avec 6 parts, titre exact, pourcentages affichés. \\
Bonus & \cel{=SI(G4>=50;"Best-seller";"")} & « Best-seller » attendu pour le switch (50), le point d'accès (93) et le câble (193). Accepter \cel{>49}. \\
\bottomrule
\end{tabularx}
\end{center}

\textbf{Répartition indicative du temps} : saisie ou ouverture 5 min, partie A 15 min, partie B 8 min, partie C 10 min,
partie D 5 min, relecture 2 min.
\fi

\end{document}
"
\documentclass[10pt,a4paper]{article}

\usepackage[T1]{fontenc}
\usepackage[utf8]{inputenc}
\usepackage{lmodern}
\usepackage[french]{babel}
\usepackage[a4paper,margin=1.5cm,top=1.3cm,bottom=1.8cm]{geometry}
\usepackage[table]{xcolor}
\usepackage{booktabs,tabularx,array}
\usepackage{enumitem}
\usepackage{fancyhdr}
\usepackage{amssymb}
\usepackage[most]{tcolorbox}
\usepackage{hyperref}
\hypersetup{hidelinks}

\providecommand{\modecorrige}{0}
\newif\ifcorrige
\ifnum\modecorrige=1 \corrigetrue \else \corrigefalse \fi

\definecolor{ipred}{HTML}{C0392B}
\definecolor{cdcbar}{HTML}{A67C3D}
\definecolor{cdcbg}{HTML}{F4F1EA}
\definecolor{qbar}{HTML}{4A6B5D}
\definecolor{qbg}{HTML}{EDF1EF}
\definecolor{corr}{HTML}{1F5FA8}
\definecolor{corrbg}{HTML}{E8F0FA}
\definecolor{grisfoot}{HTML}{828282}
\definecolor{entete}{HTML}{D9E2DD}

\newcommand{\cel}[1]{\textcolor{ipred}{\texttt{#1}}}   % référence de cellule / formule
\newcommand{\pts}[1]{\hfill\mbox{\textbf{(#1)}}}
\newcommand{\menu}[1]{\textsf{\small #1}}               % chemin de menu LibreOffice

\newtcolorbox{cdc}{enhanced,breakable,colback=cdcbg,colframe=cdcbar,boxrule=0pt,
  leftrule=3pt,arc=0pt,outer arc=0pt,left=5pt,right=5pt,top=2pt,bottom=2pt}
\newtcolorbox{partie}[2]{enhanced,breakable,colback=qbg,colframe=qbar,boxrule=0pt,
  leftrule=3pt,arc=0pt,outer arc=0pt,left=5pt,right=5pt,top=2pt,bottom=2pt,
  title={#1\hfill(#2)},coltitle=qbar,fonttitle=\bfseries,detach title,before upper={\tcbtitle\par\smallskip}}

% Zone de réponse : vierge dans le sujet, remplie dans le corrigé
\newcommand{\zone}[2]{%
  \ifcorrige
    \begin{tcolorbox}[colback=corrbg,colframe=corr,boxrule=0.6pt,arc=1pt,left=4pt,right=4pt,top=1pt,bottom=1pt,
      fontupper=\small\color{corr}]#2\end{tcolorbox}%
  \else
    \begin{tcolorbox}[colback=white,colframe=black!45,boxrule=0.5pt,arc=1pt,height=#1]\end{tcolorbox}%
  \fi}

\pagestyle{fancy}
\fancyhf{}
\renewcommand{\headrulewidth}{0pt}
\fancyfoot[L]{\small\itshape\color{grisfoot}Évaluation — Tableur LibreOffice Calc\ifcorrige\ — CORRIGÉ\fi}
\fancyfoot[R]{\small\color{grisfoot}\thepage}

\setlist{itemsep=0pt,topsep=1pt,parsep=0pt,leftmargin=*}
\setlength{\parindent}{0pt}
\setlength{\parskip}{3pt}

\begin{document}

\begin{center}
  {\Large\bfseries Évaluation pratique — Tableur (LibreOffice Calc)\par}
  \vspace{1pt}
  {\bfseries Formules, fonctions, mise en forme, tri et graphique — durée : 45 minutes\par}
  \vspace{2pt}
  BTS CIEL 1re année
  \ifcorrige\par\vspace{2pt}{\large\bfseries\color{corr}CORRIGÉ ET BARÈME — document enseignant}\fi
\end{center}

\vspace{-2pt}
{\renewcommand{\arraystretch}{1.5}
\begin{tabularx}{\linewidth}{|X|X|p{2.2cm}|}
\hline
\textbf{Nom :} & \textbf{Prénom :} & \textbf{Note :} \hfill /20 \\ \hline
\end{tabularx}}

\begin{cdc}
\textbf{Consignes.} Travail individuel, sans Internet. Documents de cours autorisés.
\begin{itemize}
  \item Ouvrir le fichier \texttt{Eval\_Tableur\_donnees.ods} (ou, à défaut, saisir le tableau ci-dessous en respectant
        \textbf{exactement} les cellules), puis l'\textbf{enregistrer immédiatement} sous
        \texttt{NOM\_Prenom\_Eval\_Tableur.ods} (\menu{Fichier > Enregistrer sous}). Enregistrer régulièrement.
  \item \textbf{Toute valeur calculée doit l'être par une formule} : un nombre recopié à la main vaut 0.
  \item Les réponses de la partie D se rédigent sur cette feuille, à rendre avec le fichier.
\end{itemize}
\end{cdc}

\textbf{Contexte.} L'entreprise \emph{NetCIEL} vend du matériel réseau. Le responsable commercial vous confie le
tableau de ses ventes 2025 (quantités vendues par trimestre) pour établir un bilan.

\begin{center}\small
\setlength{\tabcolsep}{4pt}
\begin{tabular}{|>{\columncolor{black!6}}c|l|c|c|c|c|c|c|c|c|c|}
\hline
\rowcolor{black!6} & \textbf{A} & \textbf{B} & \textbf{C} & \textbf{D} & \textbf{E} & \textbf{F} & \textbf{G} & \textbf{H} & \textbf{I} & \textbf{J} \\ \hline
1 & \multicolumn{10}{l|}{\textbf{Ventes 2025 — NetCIEL}} \\ \hline
3 & Produit & Prix unitaire & T1 & T2 & T3 & T4 & Total qté & Moy./trim. & CA annuel & Part du CA \\ \hline
4 & Switch 24 ports & 180 & 12 & 15 & 9 & 14 & & & & \\ \hline
5 & Routeur & 250 & 8 & 6 & 10 & 7 & & & & \\ \hline
6 & Point d'accès Wi-Fi & 95 & 20 & 25 & 18 & 30 & & & & \\ \hline
7 & Pare-feu & 420 & 3 & 5 & 4 & 6 & & & & \\ \hline
8 & Câble RJ45 Cat6 (lot de 10) & 35 & 40 & 55 & 38 & 60 & & & & \\ \hline
9 & Onduleur & 150 & 5 & 4 & 7 & 6 & & & & \\ \hline
10 & Total & & & & & & & & & \\ \hline
\end{tabular}
\end{center}
\vspace{-4pt}
{\footnotesize La ligne 2 reste vide. Le prix unitaire est en euros hors taxes ; CA = chiffre d'affaires.}

%---------------------------------------------------------------------
\begin{partie}{Partie A — Calculs avec formules et recopie}{8 pts}
\begin{enumerate}[label=\textbf{A\arabic*.}]
  \item En \cel{G4}, calculer le nombre total d'articles vendus dans l'année avec la fonction \cel{SOMME},
        puis recopier la formule jusqu'en \cel{G9} (poignée de recopie). \pts{1,5 pt}
  \item En \cel{H4}, calculer la moyenne trimestrielle des ventes avec la fonction \cel{MOYENNE}, puis recopier
        jusqu'en \cel{H9}. \pts{1 pt}
  \item En \cel{I4}, calculer le chiffre d'affaires annuel du produit (prix unitaire $\times$ total des quantités),
        puis recopier jusqu'en \cel{I9}. \pts{1,5 pt}
  \item Ligne 10 : calculer en \cel{C10} le total des ventes du 1\textsuperscript{er} trimestre, puis recopier
        la formule vers la droite jusqu'en \cel{G10}. Calculer aussi le CA total en \cel{I10}. \pts{1,5 pt}
  \item En \cel{J4}, calculer la part du CA du produit dans le CA total (\cel{I4} divisé par le CA total).
        La formule doit pouvoir être recopiée jusqu'en \cel{J9} \textbf{sans être modifiée} :
        utiliser une \textbf{référence absolue} (symbole \cel{\$}). \pts{1,5 pt}
  \item Saisir en \cel{A12} « CA le plus élevé », en \cel{A13} « CA le plus faible », en \cel{A14} « Écart ».
        Calculer en \cel{B12} et \cel{B13} ces valeurs à partir de la plage \cel{I4:I9} avec les fonctions
        \cel{MAX} et \cel{MIN}, puis leur écart en \cel{B14}. \pts{1 pt}
\end{enumerate}
\end{partie}

%---------------------------------------------------------------------
\begin{partie}{Partie B — Mise en forme}{4 pts}
\begin{enumerate}[label=\textbf{B\arabic*.}]
  \item Fusionner les cellules \cel{A1:J1} ; titre centré, en gras, taille 14. \pts{1 pt}
  \item En-têtes de la ligne 3 en gras, centrés, avec une couleur d'arrière-plan ; ajuster la largeur des colonnes
        pour que tout soit lisible. \pts{1 pt}
  \item Ajouter toutes les bordures au tableau \cel{A3:J10}. \pts{0,5 pt}
  \item Format \textbf{monétaire en euros} pour les colonnes \cel{B} et \cel{I} et pour \cel{B12:B14} ;
        format \textbf{pourcentage avec 1 décimale} pour la colonne \cel{J}. \pts{1,5 pt}
\end{enumerate}
\end{partie}

\newpage
%---------------------------------------------------------------------
\begin{partie}{Partie C — Tri et graphique}{5 pts}
\begin{enumerate}[label=\textbf{C\arabic*.}]
  \item Trier les produits (lignes 4 à 9 seulement, \textbf{sans la ligne Total}) par CA annuel
        \textbf{décroissant} (\menu{Données > Trier}). Vérifier que les résultats sont restés justes. \pts{1,5 pt}
  \item Insérer un graphique en \textbf{secteurs} représentant la répartition du CA annuel par produit
        (noms des produits et CA ; touche \textbf{Ctrl} pour sélectionner deux plages non adjacentes)
        (\menu{Insertion > Diagramme}). \pts{2 pts}
  \item Donner au graphique le titre « Répartition du CA 2025 » et afficher les pourcentages sur les secteurs.
        Le placer à côté du tableau, sans masquer de données. \pts{1,5 pt}
\end{enumerate}
\end{partie}

%---------------------------------------------------------------------
\begin{partie}{Partie D — Analyse (à rédiger sur cette feuille)}{3 pts}
\begin{enumerate}[label=\textbf{D\arabic*.}]
  \item Écrire la formule saisie en \cel{J4}. Que se serait-il passé en \cel{J5} si l'on avait recopié la formule
        \cel{=I4/I10}, sans le symbole \cel{\$} ? Expliquer. \pts{1 pt}
\zone{2.6cm}{\cel{=I4/\$I\$10}. Sans \cel{\$}, la recopie décale les deux références : \cel{J5} contiendrait
\cel{=I5/I11}. \cel{I11} est vide, d'où l'erreur \cel{\#DIV/0!}. Le \cel{\$} bloque la colonne et la ligne de la cellule du total.}
  \item Le produit le plus vendu en quantité est-il celui qui rapporte le plus ? Justifier avec vos résultats. \pts{1 pt}
\zone{2.4cm}{Non. Le câble RJ45 est le plus vendu (193 lots) mais ne rapporte que 6\,755~€ (15,6~\%) à cause de son
faible prix (35~€). Le switch rapporte le plus (9\,000~€, 20,8~\%) avec 50 unités seulement.}
  \item Le prix du routeur passe à 270~€. Quelles cellules sont modifiées automatiquement ? Pourquoi ? \pts{1 pt}
\zone{2.4cm}{Le CA du routeur (colonne I), le CA total \cel{I10}, toutes les parts de la colonne J, éventuellement
\cel{B12:B14}, et le graphique. Les formules utilisent des références de cellules : tout est recalculé dès qu'une donnée change.}
\end{enumerate}
\end{partie}

\textbf{Bonus (+1 pt)} — En \cel{K4}, afficher « Best-seller » si le total des quantités est supérieur ou égal
à 50, et rien sinon (fonction \cel{SI}), puis recopier jusqu'en \cel{K9}.

\textbf{En fin d'épreuve} : enregistrer, fermer le fichier, le déposer dans le dossier indiqué, rendre cette feuille.

%=====================================================================
\ifcorrige
\newpage
\section*{Corrigé — formules et résultats attendus}

LibreOffice Calc en français : noms de fonctions en français, séparateur d'arguments \cel{;}.
Valeurs ci-dessous \textbf{avant le tri} (après le tri, l'ordre des lignes change mais les valeurs restent les mêmes).

\begin{center}\small
\setlength{\tabcolsep}{4pt}
\begin{tabular}{@{}lcccc@{}}
\toprule
\textbf{Produit} & \textbf{G : Total qté} & \textbf{H : Moy./trim.} & \textbf{I : CA annuel} & \textbf{J : Part du CA} \\
\midrule
Switch 24 ports              & 50  & 12,5  & 9\,000,00~€ & 20,8~\% \\
Routeur                      & 31  & 7,75  & 7\,750,00~€ & 17,9~\% \\
Point d'accès Wi-Fi          & 93  & 23,25 & 8\,835,00~€ & 20,5~\% \\
Pare-feu                     & 18  & 4,5   & 7\,560,00~€ & 17,5~\% \\
Câble RJ45 Cat6 (lot de 10)  & 193 & 48,25 & 6\,755,00~€ & 15,6~\% \\
Onduleur                     & 22  & 5,5   & 3\,300,00~€ & 7,6~\% \\
\midrule
\textbf{Total (ligne 10)} & \textbf{407} & — & \textbf{43\,200,00~€} & (100~\%) \\
\bottomrule
\end{tabular}

\smallskip
Totaux trimestriels : \cel{C10} = 88 ; \cel{D10} = 110 ; \cel{E10} = 86 ; \cel{F10} = 123.
\quad \cel{B12} = 9\,000~€ ; \cel{B13} = 3\,300~€ ; \cel{B14} = 5\,700~€.
\end{center}

\begin{center}\small
\begin{tabularx}{\linewidth}{@{}lp{4.6cm}X@{}}
\toprule
\textbf{Question} & \textbf{Formule attendue} & \textbf{Barème et points de vigilance} \\
\midrule
A1 & \cel{=SOMME(C4:F4)} & 1 pt formule, 0,5 pt recopie. \cel{=C4+D4+E4+F4} : 1 pt sur 1,5 (fonction demandée). \\
A2 & \cel{=MOYENNE(C4:F4)} & \cel{=G4/4} accepté pour 0,5 pt sur 1 (fonction demandée). \\
A3 & \cel{=B4*G4} & 1 pt formule, 0,5 pt recopie. \\
A4 & \cel{=SOMME(C4:C9)} ; \cel{I10} : \cel{=SOMME(I4:I9)} & 1 pt ligne C–G recopiée, 0,5 pt \cel{I10}. \cel{H10} vide ou moyenne : accepté. \\
A5 & \cel{=I4/\$I\$10} & \cel{=I4/I\$10} accepté. Sans \cel{\$} : \cel{\#DIV/0!} dès \cel{J5}, 0,5 pt max. \\
A6 & \cel{=MAX(I4:I9)} ; \cel{=MIN(I4:I9)} ; \cel{=B12-B13} & 0,5 pt MAX/MIN, 0,5 pt écart. Valeur tapée : 0. \\
B1--B4 & — & Contrôle visuel. B4 : 1 pt monétaire, 0,5 pt pourcentage à 1 décimale. \\
C1 & — & Ordre attendu : Switch, Point d'accès, Routeur, Pare-feu, Câble, Onduleur. Si la ligne Total est incluse
dans le tri ou si seule la colonne I a été triée (données mélangées) : 0. \\
C2--C3 & — & Plages \cel{A4:A9} et \cel{I4:I9}. Graphique en secteurs avec 6 parts, titre exact, pourcentages affichés. \\
Bonus & \cel{=SI(G4>=50;"Best-seller";"")} & « Best-seller » attendu pour le switch (50), le point d'accès (93) et le câble (193). Accepter \cel{>49}. \\
\bottomrule
\end{tabularx}
\end{center}

\textbf{Répartition indicative du temps} : saisie ou ouverture 5 min, partie A 15 min, partie B 8 min, partie C 10 min,
partie D 5 min, relecture 2 min.
\fi

\end{document}
"
\documentclass[10pt,a4paper]{article}

\usepackage[T1]{fontenc}
\usepackage[utf8]{inputenc}
\usepackage{lmodern}
\usepackage[french]{babel}
\usepackage[a4paper,margin=1.5cm,top=1.3cm,bottom=1.8cm]{geometry}
\usepackage[table]{xcolor}
\usepackage{booktabs,tabularx,array}
\usepackage{enumitem}
\usepackage{fancyhdr}
\usepackage{amssymb}
\usepackage[most]{tcolorbox}
\usepackage{hyperref}
\hypersetup{hidelinks}

\providecommand{\modecorrige}{0}
\newif\ifcorrige
\ifnum\modecorrige=1 \corrigetrue \else \corrigefalse \fi

\definecolor{ipred}{HTML}{C0392B}
\definecolor{cdcbar}{HTML}{A67C3D}
\definecolor{cdcbg}{HTML}{F4F1EA}
\definecolor{qbar}{HTML}{4A6B5D}
\definecolor{qbg}{HTML}{EDF1EF}
\definecolor{corr}{HTML}{1F5FA8}
\definecolor{corrbg}{HTML}{E8F0FA}
\definecolor{grisfoot}{HTML}{828282}
\definecolor{entete}{HTML}{D9E2DD}

\newcommand{\cel}[1]{\textcolor{ipred}{\texttt{#1}}}   % référence de cellule / formule
\newcommand{\pts}[1]{\hfill\mbox{\textbf{(#1)}}}
\newcommand{\menu}[1]{\textsf{\small #1}}               % chemin de menu LibreOffice

\newtcolorbox{cdc}{enhanced,breakable,colback=cdcbg,colframe=cdcbar,boxrule=0pt,
  leftrule=3pt,arc=0pt,outer arc=0pt,left=5pt,right=5pt,top=2pt,bottom=2pt}
\newtcolorbox{partie}[2]{enhanced,breakable,colback=qbg,colframe=qbar,boxrule=0pt,
  leftrule=3pt,arc=0pt,outer arc=0pt,left=5pt,right=5pt,top=2pt,bottom=2pt,
  title={#1\hfill(#2)},coltitle=qbar,fonttitle=\bfseries,detach title,before upper={\tcbtitle\par\smallskip}}

% Zone de réponse : vierge dans le sujet, remplie dans le corrigé
\newcommand{\zone}[2]{%
  \ifcorrige
    \begin{tcolorbox}[colback=corrbg,colframe=corr,boxrule=0.6pt,arc=1pt,left=4pt,right=4pt,top=1pt,bottom=1pt,
      fontupper=\small\color{corr}]#2\end{tcolorbox}%
  \else
    \begin{tcolorbox}[colback=white,colframe=black!45,boxrule=0.5pt,arc=1pt,height=#1]\end{tcolorbox}%
  \fi}

\pagestyle{fancy}
\fancyhf{}
\renewcommand{\headrulewidth}{0pt}
\fancyfoot[L]{\small\itshape\color{grisfoot}Évaluation — Tableur LibreOffice Calc\ifcorrige\ — CORRIGÉ\fi}
\fancyfoot[R]{\small\color{grisfoot}\thepage}

\setlist{itemsep=0pt,topsep=1pt,parsep=0pt,leftmargin=*}
\setlength{\parindent}{0pt}
\setlength{\parskip}{3pt}

\begin{document}

\begin{center}
  {\Large\bfseries Évaluation pratique — Tableur (LibreOffice Calc)\par}
  \vspace{1pt}
  {\bfseries Formules, fonctions, mise en forme, tri et graphique — durée : 45 minutes\par}
  \vspace{2pt}
  BTS CIEL 1re année
  \ifcorrige\par\vspace{2pt}{\large\bfseries\color{corr}CORRIGÉ ET BARÈME — document enseignant}\fi
\end{center}

\vspace{-2pt}
{\renewcommand{\arraystretch}{1.5}
\begin{tabularx}{\linewidth}{|X|X|p{2.2cm}|}
\hline
\textbf{Nom :} & \textbf{Prénom :} & \textbf{Note :} \hfill /20 \\ \hline
\end{tabularx}}

\begin{cdc}
\textbf{Consignes.} Travail individuel, sans Internet. Documents de cours autorisés.
\begin{itemize}
  \item Ouvrir le fichier \texttt{Eval\_Tableur\_donnees.ods} (ou, à défaut, saisir le tableau ci-dessous en respectant
        \textbf{exactement} les cellules), puis l'\textbf{enregistrer immédiatement} sous
        \texttt{NOM\_Prenom\_Eval\_Tableur.ods} (\menu{Fichier > Enregistrer sous}). Enregistrer régulièrement.
  \item \textbf{Toute valeur calculée doit l'être par une formule} : un nombre recopié à la main vaut 0.
  \item Les réponses de la partie D se rédigent sur cette feuille, à rendre avec le fichier.
\end{itemize}
\end{cdc}

\textbf{Contexte.} L'entreprise \emph{NetCIEL} vend du matériel réseau. Le responsable commercial vous confie le
tableau de ses ventes 2025 (quantités vendues par trimestre) pour établir un bilan.

\begin{center}\small
\setlength{\tabcolsep}{4pt}
\begin{tabular}{|>{\columncolor{black!6}}c|l|c|c|c|c|c|c|c|c|c|}
\hline
\rowcolor{black!6} & \textbf{A} & \textbf{B} & \textbf{C} & \textbf{D} & \textbf{E} & \textbf{F} & \textbf{G} & \textbf{H} & \textbf{I} & \textbf{J} \\ \hline
1 & \multicolumn{10}{l|}{\textbf{Ventes 2025 — NetCIEL}} \\ \hline
3 & Produit & Prix unitaire & T1 & T2 & T3 & T4 & Total qté & Moy./trim. & CA annuel & Part du CA \\ \hline
4 & Switch 24 ports & 180 & 12 & 15 & 9 & 14 & & & & \\ \hline
5 & Routeur & 250 & 8 & 6 & 10 & 7 & & & & \\ \hline
6 & Point d'accès Wi-Fi & 95 & 20 & 25 & 18 & 30 & & & & \\ \hline
7 & Pare-feu & 420 & 3 & 5 & 4 & 6 & & & & \\ \hline
8 & Câble RJ45 Cat6 (lot de 10) & 35 & 40 & 55 & 38 & 60 & & & & \\ \hline
9 & Onduleur & 150 & 5 & 4 & 7 & 6 & & & & \\ \hline
10 & Total & & & & & & & & & \\ \hline
\end{tabular}
\end{center}
\vspace{-4pt}
{\footnotesize La ligne 2 reste vide. Le prix unitaire est en euros hors taxes ; CA = chiffre d'affaires.}

%---------------------------------------------------------------------
\begin{partie}{Partie A — Calculs avec formules et recopie}{8 pts}
\begin{enumerate}[label=\textbf{A\arabic*.}]
  \item En \cel{G4}, calculer le nombre total d'articles vendus dans l'année avec la fonction \cel{SOMME},
        puis recopier la formule jusqu'en \cel{G9} (poignée de recopie). \pts{1,5 pt}
  \item En \cel{H4}, calculer la moyenne trimestrielle des ventes avec la fonction \cel{MOYENNE}, puis recopier
        jusqu'en \cel{H9}. \pts{1 pt}
  \item En \cel{I4}, calculer le chiffre d'affaires annuel du produit (prix unitaire $\times$ total des quantités),
        puis recopier jusqu'en \cel{I9}. \pts{1,5 pt}
  \item Ligne 10 : calculer en \cel{C10} le total des ventes du 1\textsuperscript{er} trimestre, puis recopier
        la formule vers la droite jusqu'en \cel{G10}. Calculer aussi le CA total en \cel{I10}. \pts{1,5 pt}
  \item En \cel{J4}, calculer la part du CA du produit dans le CA total (\cel{I4} divisé par le CA total).
        La formule doit pouvoir être recopiée jusqu'en \cel{J9} \textbf{sans être modifiée} :
        utiliser une \textbf{référence absolue} (symbole \cel{\$}). \pts{1,5 pt}
  \item Saisir en \cel{A12} « CA le plus élevé », en \cel{A13} « CA le plus faible », en \cel{A14} « Écart ».
        Calculer en \cel{B12} et \cel{B13} ces valeurs à partir de la plage \cel{I4:I9} avec les fonctions
        \cel{MAX} et \cel{MIN}, puis leur écart en \cel{B14}. \pts{1 pt}
\end{enumerate}
\end{partie}

%---------------------------------------------------------------------
\begin{partie}{Partie B — Mise en forme}{4 pts}
\begin{enumerate}[label=\textbf{B\arabic*.}]
  \item Fusionner les cellules \cel{A1:J1} ; titre centré, en gras, taille 14. \pts{1 pt}
  \item En-têtes de la ligne 3 en gras, centrés, avec une couleur d'arrière-plan ; ajuster la largeur des colonnes
        pour que tout soit lisible. \pts{1 pt}
  \item Ajouter toutes les bordures au tableau \cel{A3:J10}. \pts{0,5 pt}
  \item Format \textbf{monétaire en euros} pour les colonnes \cel{B} et \cel{I} et pour \cel{B12:B14} ;
        format \textbf{pourcentage avec 1 décimale} pour la colonne \cel{J}. \pts{1,5 pt}
\end{enumerate}
\end{partie}

\newpage
%---------------------------------------------------------------------
\begin{partie}{Partie C — Tri et graphique}{5 pts}
\begin{enumerate}[label=\textbf{C\arabic*.}]
  \item Trier les produits (lignes 4 à 9 seulement, \textbf{sans la ligne Total}) par CA annuel
        \textbf{décroissant} (\menu{Données > Trier}). Vérifier que les résultats sont restés justes. \pts{1,5 pt}
  \item Insérer un graphique en \textbf{secteurs} représentant la répartition du CA annuel par produit
        (noms des produits et CA ; touche \textbf{Ctrl} pour sélectionner deux plages non adjacentes)
        (\menu{Insertion > Diagramme}). \pts{2 pts}
  \item Donner au graphique le titre « Répartition du CA 2025 » et afficher les pourcentages sur les secteurs.
        Le placer à côté du tableau, sans masquer de données. \pts{1,5 pt}
\end{enumerate}
\end{partie}

%---------------------------------------------------------------------
\begin{partie}{Partie D — Analyse (à rédiger sur cette feuille)}{3 pts}
\begin{enumerate}[label=\textbf{D\arabic*.}]
  \item Écrire la formule saisie en \cel{J4}. Que se serait-il passé en \cel{J5} si l'on avait recopié la formule
        \cel{=I4/I10}, sans le symbole \cel{\$} ? Expliquer. \pts{1 pt}
\zone{2.6cm}{\cel{=I4/\$I\$10}. Sans \cel{\$}, la recopie décale les deux références : \cel{J5} contiendrait
\cel{=I5/I11}. \cel{I11} est vide, d'où l'erreur \cel{\#DIV/0!}. Le \cel{\$} bloque la colonne et la ligne de la cellule du total.}
  \item Le produit le plus vendu en quantité est-il celui qui rapporte le plus ? Justifier avec vos résultats. \pts{1 pt}
\zone{2.4cm}{Non. Le câble RJ45 est le plus vendu (193 lots) mais ne rapporte que 6\,755~€ (15,6~\%) à cause de son
faible prix (35~€). Le switch rapporte le plus (9\,000~€, 20,8~\%) avec 50 unités seulement.}
  \item Le prix du routeur passe à 270~€. Quelles cellules sont modifiées automatiquement ? Pourquoi ? \pts{1 pt}
\zone{2.4cm}{Le CA du routeur (colonne I), le CA total \cel{I10}, toutes les parts de la colonne J, éventuellement
\cel{B12:B14}, et le graphique. Les formules utilisent des références de cellules : tout est recalculé dès qu'une donnée change.}
\end{enumerate}
\end{partie}

\textbf{Bonus (+1 pt)} — En \cel{K4}, afficher « Best-seller » si le total des quantités est supérieur ou égal
à 50, et rien sinon (fonction \cel{SI}), puis recopier jusqu'en \cel{K9}.

\textbf{En fin d'épreuve} : enregistrer, fermer le fichier, le déposer dans le dossier indiqué, rendre cette feuille.

%=====================================================================
\ifcorrige
\newpage
\section*{Corrigé — formules et résultats attendus}

LibreOffice Calc en français : noms de fonctions en français, séparateur d'arguments \cel{;}.
Valeurs ci-dessous \textbf{avant le tri} (après le tri, l'ordre des lignes change mais les valeurs restent les mêmes).

\begin{center}\small
\setlength{\tabcolsep}{4pt}
\begin{tabular}{@{}lcccc@{}}
\toprule
\textbf{Produit} & \textbf{G : Total qté} & \textbf{H : Moy./trim.} & \textbf{I : CA annuel} & \textbf{J : Part du CA} \\
\midrule
Switch 24 ports              & 50  & 12,5  & 9\,000,00~€ & 20,8~\% \\
Routeur                      & 31  & 7,75  & 7\,750,00~€ & 17,9~\% \\
Point d'accès Wi-Fi          & 93  & 23,25 & 8\,835,00~€ & 20,5~\% \\
Pare-feu                     & 18  & 4,5   & 7\,560,00~€ & 17,5~\% \\
Câble RJ45 Cat6 (lot de 10)  & 193 & 48,25 & 6\,755,00~€ & 15,6~\% \\
Onduleur                     & 22  & 5,5   & 3\,300,00~€ & 7,6~\% \\
\midrule
\textbf{Total (ligne 10)} & \textbf{407} & — & \textbf{43\,200,00~€} & (100~\%) \\
\bottomrule
\end{tabular}

\smallskip
Totaux trimestriels : \cel{C10} = 88 ; \cel{D10} = 110 ; \cel{E10} = 86 ; \cel{F10} = 123.
\quad \cel{B12} = 9\,000~€ ; \cel{B13} = 3\,300~€ ; \cel{B14} = 5\,700~€.
\end{center}

\begin{center}\small
\begin{tabularx}{\linewidth}{@{}lp{4.6cm}X@{}}
\toprule
\textbf{Question} & \textbf{Formule attendue} & \textbf{Barème et points de vigilance} \\
\midrule
A1 & \cel{=SOMME(C4:F4)} & 1 pt formule, 0,5 pt recopie. \cel{=C4+D4+E4+F4} : 1 pt sur 1,5 (fonction demandée). \\
A2 & \cel{=MOYENNE(C4:F4)} & \cel{=G4/4} accepté pour 0,5 pt sur 1 (fonction demandée). \\
A3 & \cel{=B4*G4} & 1 pt formule, 0,5 pt recopie. \\
A4 & \cel{=SOMME(C4:C9)} ; \cel{I10} : \cel{=SOMME(I4:I9)} & 1 pt ligne C–G recopiée, 0,5 pt \cel{I10}. \cel{H10} vide ou moyenne : accepté. \\
A5 & \cel{=I4/\$I\$10} & \cel{=I4/I\$10} accepté. Sans \cel{\$} : \cel{\#DIV/0!} dès \cel{J5}, 0,5 pt max. \\
A6 & \cel{=MAX(I4:I9)} ; \cel{=MIN(I4:I9)} ; \cel{=B12-B13} & 0,5 pt MAX/MIN, 0,5 pt écart. Valeur tapée : 0. \\
B1--B4 & — & Contrôle visuel. B4 : 1 pt monétaire, 0,5 pt pourcentage à 1 décimale. \\
C1 & — & Ordre attendu : Switch, Point d'accès, Routeur, Pare-feu, Câble, Onduleur. Si la ligne Total est incluse
dans le tri ou si seule la colonne I a été triée (données mélangées) : 0. \\
C2--C3 & — & Plages \cel{A4:A9} et \cel{I4:I9}. Graphique en secteurs avec 6 parts, titre exact, pourcentages affichés. \\
Bonus & \cel{=SI(G4>=50;"Best-seller";"")} & « Best-seller » attendu pour le switch (50), le point d'accès (93) et le câble (193). Accepter \cel{>49}. \\
\bottomrule
\end{tabularx}
\end{center}

\textbf{Répartition indicative du temps} : saisie ou ouverture 5 min, partie A 15 min, partie B 8 min, partie C 10 min,
partie D 5 min, relecture 2 min.
\fi

\end{document}
"
\documentclass[10pt,a4paper]{article}

\usepackage[T1]{fontenc}
\usepackage[utf8]{inputenc}
\usepackage{lmodern}
\usepackage[french]{babel}
\usepackage[a4paper,margin=1.5cm,top=1.3cm,bottom=1.8cm]{geometry}
\usepackage[table]{xcolor}
\usepackage{booktabs,tabularx,array}
\usepackage{enumitem}
\usepackage{fancyhdr}
\usepackage{amssymb}
\usepackage[most]{tcolorbox}
\usepackage{hyperref}
\hypersetup{hidelinks}

\providecommand{\modecorrige}{0}
\newif\ifcorrige
\ifnum\modecorrige=1 \corrigetrue \else \corrigefalse \fi

\definecolor{ipred}{HTML}{C0392B}
\definecolor{cdcbar}{HTML}{A67C3D}
\definecolor{cdcbg}{HTML}{F4F1EA}
\definecolor{qbar}{HTML}{4A6B5D}
\definecolor{qbg}{HTML}{EDF1EF}
\definecolor{corr}{HTML}{1F5FA8}
\definecolor{corrbg}{HTML}{E8F0FA}
\definecolor{grisfoot}{HTML}{828282}
\definecolor{entete}{HTML}{D9E2DD}

\newcommand{\cel}[1]{\textcolor{ipred}{\texttt{#1}}}   % référence de cellule / formule
\newcommand{\pts}[1]{\hfill\mbox{\textbf{(#1)}}}
\newcommand{\menu}[1]{\textsf{\small #1}}               % chemin de menu LibreOffice

\newtcolorbox{cdc}{enhanced,breakable,colback=cdcbg,colframe=cdcbar,boxrule=0pt,
  leftrule=3pt,arc=0pt,outer arc=0pt,left=5pt,right=5pt,top=2pt,bottom=2pt}
\newtcolorbox{partie}[2]{enhanced,breakable,colback=qbg,colframe=qbar,boxrule=0pt,
  leftrule=3pt,arc=0pt,outer arc=0pt,left=5pt,right=5pt,top=2pt,bottom=2pt,
  title={#1\hfill(#2)},coltitle=qbar,fonttitle=\bfseries,detach title,before upper={\tcbtitle\par\smallskip}}

% Zone de réponse : vierge dans le sujet, remplie dans le corrigé
\newcommand{\zone}[2]{%
  \ifcorrige
    \begin{tcolorbox}[colback=corrbg,colframe=corr,boxrule=0.6pt,arc=1pt,left=4pt,right=4pt,top=1pt,bottom=1pt,
      fontupper=\small\color{corr}]#2\end{tcolorbox}%
  \else
    \begin{tcolorbox}[colback=white,colframe=black!45,boxrule=0.5pt,arc=1pt,height=#1]\end{tcolorbox}%
  \fi}

\pagestyle{fancy}
\fancyhf{}
\renewcommand{\headrulewidth}{0pt}
\fancyfoot[L]{\small\itshape\color{grisfoot}Évaluation — Tableur LibreOffice Calc\ifcorrige\ — CORRIGÉ\fi}
\fancyfoot[R]{\small\color{grisfoot}\thepage}

\setlist{itemsep=0pt,topsep=1pt,parsep=0pt,leftmargin=*}
\setlength{\parindent}{0pt}
\setlength{\parskip}{3pt}

\begin{document}

\begin{center}
  {\Large\bfseries Évaluation pratique — Tableur (LibreOffice Calc)\par}
  \vspace{1pt}
  {\bfseries Formules, fonctions, mise en forme, tri et graphique — durée : 45 minutes\par}
  \vspace{2pt}
  BTS CIEL 1re année
  \ifcorrige\par\vspace{2pt}{\large\bfseries\color{corr}CORRIGÉ ET BARÈME — document enseignant}\fi
\end{center}

\vspace{-2pt}
{\renewcommand{\arraystretch}{1.5}
\begin{tabularx}{\linewidth}{|X|X|p{2.2cm}|}
\hline
\textbf{Nom :} & \textbf{Prénom :} & \textbf{Note :} \hfill /20 \\ \hline
\end{tabularx}}

\begin{cdc}
\textbf{Consignes.} Travail individuel, sans Internet. Documents de cours autorisés.
\begin{itemize}
  \item Ouvrir le fichier \texttt{Eval\_Tableur\_donnees.ods} (ou, à défaut, saisir le tableau ci-dessous en respectant
        \textbf{exactement} les cellules), puis l'\textbf{enregistrer immédiatement} sous
        \texttt{NOM\_Prenom\_Eval\_Tableur.ods} (\menu{Fichier > Enregistrer sous}). Enregistrer régulièrement.
  \item \textbf{Toute valeur calculée doit l'être par une formule} : un nombre recopié à la main vaut 0.
  \item Les réponses de la partie D se rédigent sur cette feuille, à rendre avec le fichier.
\end{itemize}
\end{cdc}

\textbf{Contexte.} L'entreprise \emph{NetCIEL} vend du matériel réseau. Le responsable commercial vous confie le
tableau de ses ventes 2025 (quantités vendues par trimestre) pour établir un bilan.

\begin{center}\small
\setlength{\tabcolsep}{4pt}
\begin{tabular}{|>{\columncolor{black!6}}c|l|c|c|c|c|c|c|c|c|c|}
\hline
\rowcolor{black!6} & \textbf{A} & \textbf{B} & \textbf{C} & \textbf{D} & \textbf{E} & \textbf{F} & \textbf{G} & \textbf{H} & \textbf{I} & \textbf{J} \\ \hline
1 & \multicolumn{10}{l|}{\textbf{Ventes 2025 — NetCIEL}} \\ \hline
3 & Produit & Prix unitaire & T1 & T2 & T3 & T4 & Total qté & Moy./trim. & CA annuel & Part du CA \\ \hline
4 & Switch 24 ports & 180 & 12 & 15 & 9 & 14 & & & & \\ \hline
5 & Routeur & 250 & 8 & 6 & 10 & 7 & & & & \\ \hline
6 & Point d'accès Wi-Fi & 95 & 20 & 25 & 18 & 30 & & & & \\ \hline
7 & Pare-feu & 420 & 3 & 5 & 4 & 6 & & & & \\ \hline
8 & Câble RJ45 Cat6 (lot de 10) & 35 & 40 & 55 & 38 & 60 & & & & \\ \hline
9 & Onduleur & 150 & 5 & 4 & 7 & 6 & & & & \\ \hline
10 & Total & & & & & & & & & \\ \hline
\end{tabular}
\end{center}
\vspace{-4pt}
{\footnotesize La ligne 2 reste vide. Le prix unitaire est en euros hors taxes ; CA = chiffre d'affaires.}

%---------------------------------------------------------------------
\begin{partie}{Partie A — Calculs avec formules et recopie}{8 pts}
\begin{enumerate}[label=\textbf{A\arabic*.}]
  \item En \cel{G4}, calculer le nombre total d'articles vendus dans l'année avec la fonction \cel{SOMME},
        puis recopier la formule jusqu'en \cel{G9} (poignée de recopie). \pts{1,5 pt}
  \item En \cel{H4}, calculer la moyenne trimestrielle des ventes avec la fonction \cel{MOYENNE}, puis recopier
        jusqu'en \cel{H9}. \pts{1 pt}
  \item En \cel{I4}, calculer le chiffre d'affaires annuel du produit (prix unitaire $\times$ total des quantités),
        puis recopier jusqu'en \cel{I9}. \pts{1,5 pt}
  \item Ligne 10 : calculer en \cel{C10} le total des ventes du 1\textsuperscript{er} trimestre, puis recopier
        la formule vers la droite jusqu'en \cel{G10}. Calculer aussi le CA total en \cel{I10}. \pts{1,5 pt}
  \item En \cel{J4}, calculer la part du CA du produit dans le CA total (\cel{I4} divisé par le CA total).
        La formule doit pouvoir être recopiée jusqu'en \cel{J9} \textbf{sans être modifiée} :
        utiliser une \textbf{référence absolue} (symbole \cel{\$}). \pts{1,5 pt}
  \item Saisir en \cel{A12} « CA le plus élevé », en \cel{A13} « CA le plus faible », en \cel{A14} « Écart ».
        Calculer en \cel{B12} et \cel{B13} ces valeurs à partir de la plage \cel{I4:I9} avec les fonctions
        \cel{MAX} et \cel{MIN}, puis leur écart en \cel{B14}. \pts{1 pt}
\end{enumerate}
\end{partie}

%---------------------------------------------------------------------
\begin{partie}{Partie B — Mise en forme}{4 pts}
\begin{enumerate}[label=\textbf{B\arabic*.}]
  \item Fusionner les cellules \cel{A1:J1} ; titre centré, en gras, taille 14. \pts{1 pt}
  \item En-têtes de la ligne 3 en gras, centrés, avec une couleur d'arrière-plan ; ajuster la largeur des colonnes
        pour que tout soit lisible. \pts{1 pt}
  \item Ajouter toutes les bordures au tableau \cel{A3:J10}. \pts{0,5 pt}
  \item Format \textbf{monétaire en euros} pour les colonnes \cel{B} et \cel{I} et pour \cel{B12:B14} ;
        format \textbf{pourcentage avec 1 décimale} pour la colonne \cel{J}. \pts{1,5 pt}
\end{enumerate}
\end{partie}

\newpage
%---------------------------------------------------------------------
\begin{partie}{Partie C — Tri et graphique}{5 pts}
\begin{enumerate}[label=\textbf{C\arabic*.}]
  \item Trier les produits (lignes 4 à 9 seulement, \textbf{sans la ligne Total}) par CA annuel
        \textbf{décroissant} (\menu{Données > Trier}). Vérifier que les résultats sont restés justes. \pts{1,5 pt}
  \item Insérer un graphique en \textbf{secteurs} représentant la répartition du CA annuel par produit
        (noms des produits et CA ; touche \textbf{Ctrl} pour sélectionner deux plages non adjacentes)
        (\menu{Insertion > Diagramme}). \pts{2 pts}
  \item Donner au graphique le titre « Répartition du CA 2025 » et afficher les pourcentages sur les secteurs.
        Le placer à côté du tableau, sans masquer de données. \pts{1,5 pt}
\end{enumerate}
\end{partie}

%---------------------------------------------------------------------
\begin{partie}{Partie D — Analyse (à rédiger sur cette feuille)}{3 pts}
\begin{enumerate}[label=\textbf{D\arabic*.}]
  \item Écrire la formule saisie en \cel{J4}. Que se serait-il passé en \cel{J5} si l'on avait recopié la formule
        \cel{=I4/I10}, sans le symbole \cel{\$} ? Expliquer. \pts{1 pt}
\zone{2.6cm}{\cel{=I4/\$I\$10}. Sans \cel{\$}, la recopie décale les deux références : \cel{J5} contiendrait
\cel{=I5/I11}. \cel{I11} est vide, d'où l'erreur \cel{\#DIV/0!}. Le \cel{\$} bloque la colonne et la ligne de la cellule du total.}
  \item Le produit le plus vendu en quantité est-il celui qui rapporte le plus ? Justifier avec vos résultats. \pts{1 pt}
\zone{2.4cm}{Non. Le câble RJ45 est le plus vendu (193 lots) mais ne rapporte que 6\,755~€ (15,6~\%) à cause de son
faible prix (35~€). Le switch rapporte le plus (9\,000~€, 20,8~\%) avec 50 unités seulement.}
  \item Le prix du routeur passe à 270~€. Quelles cellules sont modifiées automatiquement ? Pourquoi ? \pts{1 pt}
\zone{2.4cm}{Le CA du routeur (colonne I), le CA total \cel{I10}, toutes les parts de la colonne J, éventuellement
\cel{B12:B14}, et le graphique. Les formules utilisent des références de cellules : tout est recalculé dès qu'une donnée change.}
\end{enumerate}
\end{partie}

\textbf{Bonus (+1 pt)} — En \cel{K4}, afficher « Best-seller » si le total des quantités est supérieur ou égal
à 50, et rien sinon (fonction \cel{SI}), puis recopier jusqu'en \cel{K9}.

\textbf{En fin d'épreuve} : enregistrer, fermer le fichier, le déposer dans le dossier indiqué, rendre cette feuille.

%=====================================================================
\ifcorrige
\newpage
\section*{Corrigé — formules et résultats attendus}

LibreOffice Calc en français : noms de fonctions en français, séparateur d'arguments \cel{;}.
Valeurs ci-dessous \textbf{avant le tri} (après le tri, l'ordre des lignes change mais les valeurs restent les mêmes).

\begin{center}\small
\setlength{\tabcolsep}{4pt}
\begin{tabular}{@{}lcccc@{}}
\toprule
\textbf{Produit} & \textbf{G : Total qté} & \textbf{H : Moy./trim.} & \textbf{I : CA annuel} & \textbf{J : Part du CA} \\
\midrule
Switch 24 ports              & 50  & 12,5  & 9\,000,00~€ & 20,8~\% \\
Routeur                      & 31  & 7,75  & 7\,750,00~€ & 17,9~\% \\
Point d'accès Wi-Fi          & 93  & 23,25 & 8\,835,00~€ & 20,5~\% \\
Pare-feu                     & 18  & 4,5   & 7\,560,00~€ & 17,5~\% \\
Câble RJ45 Cat6 (lot de 10)  & 193 & 48,25 & 6\,755,00~€ & 15,6~\% \\
Onduleur                     & 22  & 5,5   & 3\,300,00~€ & 7,6~\% \\
\midrule
\textbf{Total (ligne 10)} & \textbf{407} & — & \textbf{43\,200,00~€} & (100~\%) \\
\bottomrule
\end{tabular}

\smallskip
Totaux trimestriels : \cel{C10} = 88 ; \cel{D10} = 110 ; \cel{E10} = 86 ; \cel{F10} = 123.
\quad \cel{B12} = 9\,000~€ ; \cel{B13} = 3\,300~€ ; \cel{B14} = 5\,700~€.
\end{center}

\begin{center}\small
\begin{tabularx}{\linewidth}{@{}lp{4.6cm}X@{}}
\toprule
\textbf{Question} & \textbf{Formule attendue} & \textbf{Barème et points de vigilance} \\
\midrule
A1 & \cel{=SOMME(C4:F4)} & 1 pt formule, 0,5 pt recopie. \cel{=C4+D4+E4+F4} : 1 pt sur 1,5 (fonction demandée). \\
A2 & \cel{=MOYENNE(C4:F4)} & \cel{=G4/4} accepté pour 0,5 pt sur 1 (fonction demandée). \\
A3 & \cel{=B4*G4} & 1 pt formule, 0,5 pt recopie. \\
A4 & \cel{=SOMME(C4:C9)} ; \cel{I10} : \cel{=SOMME(I4:I9)} & 1 pt ligne C–G recopiée, 0,5 pt \cel{I10}. \cel{H10} vide ou moyenne : accepté. \\
A5 & \cel{=I4/\$I\$10} & \cel{=I4/I\$10} accepté. Sans \cel{\$} : \cel{\#DIV/0!} dès \cel{J5}, 0,5 pt max. \\
A6 & \cel{=MAX(I4:I9)} ; \cel{=MIN(I4:I9)} ; \cel{=B12-B13} & 0,5 pt MAX/MIN, 0,5 pt écart. Valeur tapée : 0. \\
B1--B4 & — & Contrôle visuel. B4 : 1 pt monétaire, 0,5 pt pourcentage à 1 décimale. \\
C1 & — & Ordre attendu : Switch, Point d'accès, Routeur, Pare-feu, Câble, Onduleur. Si la ligne Total est incluse
dans le tri ou si seule la colonne I a été triée (données mélangées) : 0. \\
C2--C3 & — & Plages \cel{A4:A9} et \cel{I4:I9}. Graphique en secteurs avec 6 parts, titre exact, pourcentages affichés. \\
Bonus & \cel{=SI(G4>=50;"Best-seller";"")} & « Best-seller » attendu pour le switch (50), le point d'accès (93) et le câble (193). Accepter \cel{>49}. \\
\bottomrule
\end{tabularx}
\end{center}

\textbf{Répartition indicative du temps} : saisie ou ouverture 5 min, partie A 15 min, partie B 8 min, partie C 10 min,
partie D 5 min, relecture 2 min.
\fi

\end{document}
